\documentclass[11pt, a4paper]{article}

% bfw-style.tex — gemeinsame Style-Definitionen für alle Handouts/Aufgabenhefte
% Voraussetzung: Kompilation mit xelatex (wegen fontspec)

\usepackage[ngerman]{babel}
\usepackage{geometry}
\geometry{a4paper, top=2.2cm, bottom=2.2cm, left=2.3cm, right=2.3cm}

\usepackage{fontspec}
% Fallback-Kette: Source Sans Pro -> Calibri -> Segoe UI -> Latin Modern Sans
% (Latin Modern ist immer mit MiKTeX vorhanden)
\IfFontExistsTF{Source Sans Pro}{%
  \setmainfont{Source Sans Pro}[Ligatures=TeX]%
}{\IfFontExistsTF{Calibri}{%
  \setmainfont{Calibri}[Ligatures=TeX]%
}{\IfFontExistsTF{Segoe UI}{%
  \setmainfont{Segoe UI}[Ligatures=TeX]%
}{%
  \setmainfont{Latin Modern Sans}[Ligatures=TeX]%
}}}
\IfFontExistsTF{Source Code Pro}{%
  \setmonofont{Source Code Pro}[Scale=0.92]%
}{\IfFontExistsTF{Consolas}{%
  \setmonofont{Consolas}[Scale=0.92]%
}{%
  \setmonofont{Latin Modern Mono}[Scale=0.92]%
}}

\usepackage{xcolor}
% Farbpalette: hell, freundlich, modern
\definecolor{bfwPrimary}{HTML}{0EA5A4}    % Petrol/Türkis - Primär
\definecolor{bfwPrimaryDark}{HTML}{0F766E} % dunkler Petrol für Headings
\definecolor{bfwAccent}{HTML}{F59E0B}     % Amber - Akzent
\definecolor{bfwSuccess}{HTML}{10B981}    % Grün - Erfolg / Lösung
\definecolor{bfwWarn}{HTML}{F97316}       % Orange - Achtung
\definecolor{bfwInk}{HTML}{0F172A}        % fast Schwarz
\definecolor{bfwInkSoft}{HTML}{475569}    % gedämpftes Grau
\definecolor{bfwBgSoft}{HTML}{F8FAFC}     % off-white
\definecolor{bfwBgTeal}{HTML}{ECFEFF}     % helles Türkis
\definecolor{bfwBgAmber}{HTML}{FEF3C7}    % helles Gelb
\definecolor{bfwBgRose}{HTML}{FFE4E6}     % helles Rosé für Eselsbrücken

\usepackage[colorlinks=true, linkcolor=bfwPrimaryDark, urlcolor=bfwPrimaryDark]{hyperref}
\usepackage{enumitem}
\setlist[itemize]{leftmargin=*, itemsep=2pt, topsep=2pt}
\setlist[enumerate]{leftmargin=*, itemsep=2pt, topsep=2pt}

\usepackage[most]{tcolorbox}
\usepackage{booktabs}
\usepackage{tikz}
\usetikzlibrary{decorations.pathmorphing, positioning, shapes.geometric, arrows.meta}

% Merkkästchen — wichtige Definition / Kernpunkt
\newtcolorbox{merksatz}[1][]{%
  enhanced, breakable,
  colback=bfwBgTeal,
  colframe=bfwPrimary,
  boxrule=0.6pt,
  arc=4pt,
  left=10pt, right=10pt, top=8pt, bottom=8pt,
  fonttitle=\bfseries\color{bfwPrimaryDark},
  title={\faLightbulb\ Merksatz},
  #1
}

% Eselsbrücke — Mnemoniks
\newtcolorbox{eselsbruecke}[1][]{%
  enhanced, breakable,
  colback=bfwBgRose,
  colframe=bfwAccent,
  boxrule=0.6pt,
  arc=4pt,
  left=10pt, right=10pt, top=8pt, bottom=8pt,
  fonttitle=\bfseries\color{bfwWarn},
  title={Eselsbrücke},
  #1
}

% Beispiel-Box
\newtcolorbox{beispiel}[1][]{%
  enhanced, breakable,
  colback=bfwBgSoft,
  colframe=bfwInkSoft,
  boxrule=0.4pt,
  arc=3pt,
  left=10pt, right=10pt, top=8pt, bottom=8pt,
  fonttitle=\bfseries\color{bfwInk},
  title={Beispiel},
  #1
}

% Lösungs-Box (für Aufgabenhefte)
\newtcolorbox{loesung}[1][]{%
  enhanced, breakable,
  colback=bfwBgAmber!40,
  colframe=bfwSuccess,
  boxrule=0.6pt,
  arc=3pt,
  left=10pt, right=10pt, top=8pt, bottom=8pt,
  fonttitle=\bfseries\color{bfwSuccess!70!black},
  title={Lösung},
  #1
}

% Glossar-Eintrag
\newcommand{\glossareintrag}[2]{%
  \par\noindent\textbf{\color{bfwPrimaryDark}#1} \enspace #2\par\smallskip
}

% Section-Stil — etwas farbiger als Standard
\usepackage{titlesec}
\titleformat{\section}
  {\Large\bfseries\color{bfwPrimaryDark}}
  {\thesection}{0.6em}{}
\titleformat{\subsection}
  {\large\bfseries\color{bfwPrimary}}
  {\thesubsection}{0.5em}{}
\titleformat{\subsubsection}
  {\normalsize\bfseries\color{bfwInk}}
  {\thesubsubsection}{0.4em}{}

% TikZ-Stil "skizzenhaft" — etwas weicher als Rough.js, aber stabil
\tikzset{
  roughbox/.style = {
    draw=bfwPrimaryDark, thick, fill=bfwBgTeal,
    rounded corners=4pt, inner sep=6pt
  },
  roughaccent/.style = {
    draw=bfwAccent, thick, fill=bfwBgAmber,
    rounded corners=4pt, inner sep=6pt
  },
  roughline/.style = {
    draw=bfwInkSoft, thick, -{Stealth[length=2.5mm]}
  }
}

% Bullet-Symbole (bewusst kein FontAwesome — vermeidet Font-Installations-Probleme)
\newcommand{\faLightbulb}{$\bullet$}
\newcommand{\faBookmark}{$\bullet$}

% Header/Footer
\usepackage{fancyhdr}
\pagestyle{fancy}
\fancyhf{}
\renewcommand{\headrulewidth}{0pt}
\fancyfoot[L]{\small\color{bfwInkSoft}\thepage}
\fancyfoot[R]{\small\color{bfwInkSoft} BFW M-064 Beschaffung im Büromanagement}


\title{\vspace*{-1.5cm}\Huge\bfseries\color{bfwPrimaryDark}Aufgabenheft Session~1\\[0.4em]
  \large\color{bfwInkSoft}\textnormal{Rechtliche Grundlagen der Berufsausbildung --- Übungen im IHK-Klausurformat}}
\author{\textsf{Lernfeld 1 --- Die eigene Rolle im Betrieb mitgestalten \textperiodcentered{} \small\copyright\ Can Siebert\,\textperiodcentered\,2026}}
\date{Selbstlernmaterial}

\begin{document}
\maketitle
\thispagestyle{empty}

\begin{merksatz}[title={\bfseries So nutzen Sie dieses Aufgabenheft}]
Bearbeiten Sie alle Aufgaben zunächst \emph{ohne} in die Lösungen zu schauen. Schreiben
Sie Ihre Antworten in vollständigen Sätzen --- die IHK-Prüfung verlangt formulierte
Begründungen, nicht bloße Stichworte. Beachten Sie die \emph{Operatoren} (nennen,
erläutern, beurteilen, begründen) --- sie geben den geforderten Antwort-Tiefegrad vor.
Erst danach öffnen Sie den Lösungsteil ab Seite~\pageref{loesungen} und vergleichen.
\end{merksatz}

\tableofcontents
\vspace{1em}
\hrule

\section{Aufgaben}

\subsection*{Aufgabe~1 --- Duales System und Rollen (10 Punkte)}

\noindent\textbf{\color{bfwAccent}StudyFlix-Vorbereitung:}\enspace \href{https://studyflix.de/ausbildung/karriere-ausbildung/duale-ausbildung-4787}{StudyFlix: Duale Ausbildung}

\medskip
Sie beginnen Ihre Umschulung zur Kauffrau/zum Kaufmann für Büromanagement.

\begin{enumerate}[label=(\alph*)]
  \item \textbf{Erklären} Sie, warum die Berufsausbildung als „duale Ausbildung" bezeichnet wird. \hfill (3~P.)
  \item \textbf{Nennen} Sie die beiden Lernorte und je eine typische Aufgabe. \hfill (3~P.)
  \item \textbf{Unterscheiden} Sie die Begriffe Ausbildende:r, Ausbilder:in und Auszubildende:r. \hfill (4~P.)
\end{enumerate}

\subsection*{Aufgabe~2 --- Rechtsgrundlagen (12 Punkte)}

\noindent\textbf{\color{bfwAccent}StudyFlix-Vorbereitung:}\enspace \href{https://studyflix.de/ausbildung/karriere-ausbildung/ausbildungsvertrag-4788}{StudyFlix: Ausbildungsvertrag (BBiG, Ausbildungsordnung)}

\medskip
Die Rechtsgrundlagen der Ausbildung bauen vom Allgemeinen zum Konkreten aufeinander auf.

\begin{enumerate}[label=(\alph*)]
  \item \textbf{Beschreiben} Sie, was das BBiG grundsätzlich regelt. \hfill (3~P.)
  \item \textbf{Erläutern} Sie den Unterschied zwischen Ausbildungsrahmenplan und betrieblichem Ausbildungsplan. \hfill (4~P.)
  \item \textbf{Ordnen} Sie in die richtige Reihenfolge (vom Allgemeinen zum Konkreten): betrieblicher Ausbildungsplan, BBiG, Ausbildungsrahmenplan, Ausbildungsordnung. \hfill (2~P.)
  \item \textbf{Erklären} Sie, warum der betriebliche Ausbildungsplan für jeden Betrieb individuell ist. \hfill (3~P.)
\end{enumerate}

\subsection*{Aufgabe~3 --- Der Berufsausbildungsvertrag (14 Punkte)}

\noindent\textbf{\color{bfwAccent}StudyFlix-Vorbereitung:}\enspace \href{https://studyflix.de/ausbildung/karriere-ausbildung/ausbildungsvertrag-4788}{StudyFlix: Ausbildungsvertrag --- § 11 BBiG}

\medskip
Ein 17-jähriger Bewerber soll einen Ausbildungsvertrag erhalten.

\begin{enumerate}[label=(\alph*)]
  \item \textbf{Nennen} Sie, wer den Vertrag unterschreiben muss. \hfill (3~P.)
  \item \textbf{Erklären} Sie, in welcher Form der wesentliche Vertragsinhalt niederzulegen ist und welche Form ausgeschlossen ist. \hfill (3~P.)
  \item \textbf{Nennen} Sie fünf der neun Pflichtinhalte nach \S~11 BBiG. \hfill (5~P.)
  \item \textbf{Erläutern} Sie, bei welcher Stelle das Ausbildungsverhältnis eingetragen wird und warum. \hfill (3~P.)
\end{enumerate}

\subsection*{Aufgabe~4 --- Pflichten und Fristen (12 Punkte)}

\noindent\textbf{\color{bfwAccent}StudyFlix-Vorbereitung:}\enspace \href{https://studyflix.de/ausbildung/karriere-ausbildung/rechte-und-pflichten-als-azubi-4792}{StudyFlix: Rechte und Pflichten als Azubi}

\begin{enumerate}[label=(\alph*)]
  \item \textbf{Nennen} Sie drei Pflichten der Auszubildenden (\S~13) und drei Pflichten der Ausbildenden (\S\S~14--16). \hfill (6~P.)
  \item \textbf{Erläutern} Sie die zulässige Dauer der Probezeit (\S~20 BBiG). \hfill (2~P.)
  \item \textbf{Erklären} Sie, warum die Ausbildungsvergütung mindestens jährlich steigen muss (\S~17 BBiG). \hfill (2~P.)
  \item \textbf{Ordnen} Sie zu: Wer hat die Pflicht zur Freistellung für die Berufsschule und wer zum Führen des Berichtshefts? \hfill (2~P.)
\end{enumerate}

\subsection*{Aufgabe~5 --- Jugendarbeitsschutzgesetz (14 Punkte)}

\noindent\textbf{\color{bfwAccent}StudyFlix-Vorbereitung:}\enspace \href{https://studyflix.de/suche?q=Jugendarbeitsschutzgesetz}{StudyFlix: Jugendarbeitsschutzgesetz (Suche)}

\medskip
Lena (17) beginnt ihre Ausbildung. Prüfen Sie mithilfe des JArbSchG:

\begin{enumerate}[label=(\alph*)]
  \item \textbf{Beurteilen} Sie: Der Betrieb möchte Lena 9~Stunden täglich an 6~Tagen beschäftigen. Zulässig? \hfill (4~P.)
  \item \textbf{Nennen} Sie die zulässige Beschäftigungszeit (Uhrzeiten) und die Mindest-Nachtruhe. \hfill (3~P.)
  \item \textbf{Nennen} Sie Lenas Mindesturlaub und die vor Beginn nötige ärztliche Untersuchung. \hfill (4~P.)
  \item \textbf{Nennen} Sie zwei Tätigkeiten, die für Jugendliche verboten sind. \hfill (3~P.)
\end{enumerate}

\subsection*{Aufgabe~6 --- Kündigung, Zeugnis und Mitbestimmung (14 Punkte)}

\noindent\textbf{\color{bfwAccent}StudyFlix-Vorbereitung:}\enspace \href{https://studyflix.de/ausbildung/karriere-ausbildung/kuendigung-ausbildung-4804}{StudyFlix: Kündigung Ausbildung (§ 22 BBiG)}

\begin{enumerate}[label=(\alph*)]
  \item \textbf{Erläutern} Sie die Kündigungsmöglichkeiten während und nach der Probezeit (\S~22 BBiG). \hfill (5~P.)
  \item \textbf{Unterscheiden} Sie einfaches und qualifiziertes Ausbildungszeugnis (\S~16 BBiG). \hfill (4~P.)
  \item \textbf{Nennen} Sie, ab wie vielen Arbeitnehmern ein Betriebsrat gewählt wird und wie lange seine Amtszeit dauert. \hfill (2~P.)
  \item \textbf{Beschreiben} Sie kurz eine Aufgabe der Jugend- und Auszubildendenvertretung (JAV). \hfill (3~P.)
\end{enumerate}

\vspace{1em}
\noindent\textbf{Punkte gesamt: 76}

\newpage

\section{Lösungen}\label{loesungen}

\subsection*{Lösung Aufgabe~1}
\begin{loesung}
\textbf{(a)} Die Ausbildung findet parallel an \emph{zwei} Lernorten statt --- deshalb
„dual": im Betrieb (praktische Ausbildung) und in der Berufsschule (theoretische
Ausbildung).

\textbf{(b)} \emph{Betrieb}: Vermittlung der praktischen beruflichen Handlungsfähigkeit.
\emph{Berufsschule}: Vermittlung fachtheoretischer und allgemeinbildender Kenntnisse.

\textbf{(c)} \emph{Ausbildende:r} = Vertragspartner, das Unternehmen/die Inhaberin/der
Inhaber, das andere zur Ausbildung einstellt. \emph{Ausbilder:in} = die Person, die die
Ausbildung im Betrieb praktisch durchführt und dazu beauftragt ist. \emph{Auszubildende:r}
= die Person, die ausgebildet wird.
\end{loesung}

\subsection*{Lösung Aufgabe~2}
\begin{loesung}
\textbf{(a)} Das BBiG (in Kraft seit 1969, novelliert mit Wirkung ab 01.01.2020) regelt in
Deutschland die Berufsausbildung: u.\,a.\ Vertrag, Pflichten, Vergütung, Prüfungen und
Beendigung.

\textbf{(b)} Der \emph{Ausbildungsrahmenplan} ist Teil der Ausbildungsordnung und gibt eine
grobe, bundeseinheitliche sachliche und zeitliche Gliederung vor. Der \emph{betriebliche
Ausbildungsplan} überträgt diese Vorgaben konkret auf das jeweilige Unternehmen und ist
Anlage des Ausbildungsvertrages.

\textbf{(c)} Reihenfolge vom Allgemeinen zum Konkreten: \textbf{BBiG $\rightarrow$
Ausbildungsordnung $\rightarrow$ Ausbildungsrahmenplan $\rightarrow$ betrieblicher
Ausbildungsplan}.

\textbf{(d)} Weil jeder Betrieb andere Abteilungen, Aufgaben und zeitliche Abläufe hat; der
Plan muss auf das jeweilige Unternehmen abgestimmt sein, damit alle vorgeschriebenen
Inhalte tatsächlich vermittelt werden können.
\end{loesung}

\subsection*{Lösung Aufgabe~3}
\begin{loesung}
\textbf{(a)} Die/der Ausbildende und die/der Auszubildende; da der Bewerber unter 18 ist,
zusätzlich die gesetzlichen Vertreter (i.\,d.\,R.\ beide Elternteile).

\textbf{(b)} Der wesentliche Vertragsinhalt ist \emph{schriftlich} niederzulegen (\S~11
Abs.~1 BBiG); die \emph{elektronische Form ist ausgeschlossen}.

\textbf{(c)} Fünf von neun (Beispiele): Art, sachliche/zeitliche Gliederung und Ziel;
Beginn und Dauer; Ausbildungsmaßnahmen außerhalb der Ausbildungsstätte; tägliche
Ausbildungszeit; Probezeit; Vergütung; Urlaub; Kündigungsvoraussetzungen; Hinweis auf
Tarifverträge/Betriebsvereinbarungen.

\textbf{(d)} Bei der \emph{zuständigen Stelle}, für kaufmännische Berufe die \emph{IHK}. Sie
führt das Verzeichnis der Berufsausbildungsverhältnisse und überwacht die Ausbildung;
dadurch wird das Verhältnis offiziell registriert und geschützt.
\end{loesung}

\subsection*{Lösung Aufgabe~4}
\begin{loesung}
\textbf{(a)} \emph{Auszubildende (\S~13)}: Aufgaben sorgfältig ausführen; Weisungen folgen;
Betriebsordnung beachten; Werkzeug pfleglich behandeln; Stillschweigen wahren; Berichtsheft
führen. \emph{Ausbildende (\S\S~14--16)}: Handlungsfähigkeit vermitteln; Ausbildungsmittel
kostenlos stellen; zum Berufsschulbesuch anhalten und freistellen; Fürsorge; Vergütung
zahlen; Zeugnis ausstellen.

\textbf{(b)} Die Probezeit muss mindestens \emph{einen Monat} und darf höchstens \emph{vier
Monate} betragen (\S~20 BBiG).

\textbf{(c)} Weil \S~17 BBiG dies vorschreibt: Mit zunehmender Ausbildungsdauer wächst die
Leistungsfähigkeit der Auszubildenden; die Vergütung soll dem folgen und muss daher
mindestens jährlich steigen (angemessene Vergütung).

\textbf{(d)} Die \emph{Freistellung} für die Berufsschule ist Pflicht der/des Ausbildenden
(\S~15). Das \emph{Führen des Berichtshefts} ist Pflicht der/des Auszubildenden (\S~13); die
Ausbildenden müssen dazu anhalten und den Nachweis durchsehen (\S~14).
\end{loesung}

\subsection*{Lösung Aufgabe~5}
\begin{loesung}
\textbf{(a)} Nicht zulässig. Jugendliche (unter 18) dürfen in der Regel höchstens
\emph{8~Stunden täglich} und \emph{40~Stunden wöchentlich} sowie nur an \emph{fünf} Tagen
pro Woche beschäftigt werden (\S\S~8, 15 JArbSchG). 9~Stunden an 6~Tagen verstößt gleich
mehrfach dagegen.

\textbf{(b)} Beschäftigung grundsätzlich nur von \emph{6 bis 20~Uhr} (\S~14); nach
Arbeitsende ist eine ununterbrochene Freizeit/Nachtruhe von mindestens \emph{12~Stunden}
einzuhalten (\S~13).

\textbf{(c)} Als noch nicht 18-Jährige mindestens \emph{25~Werktage} Urlaub (gestaffelt nach
\S~19: unter 16 mind.\ 30, unter 17 mind.\ 27, unter 18 mind.\ 25). Vor Ausbildungsbeginn
muss eine \emph{Erstuntersuchung} vorliegen (\S~32); innerhalb des ersten Jahres folgt die
Nachuntersuchung (\S~33).

\textbf{(d)} Zwei von: gefährliche Arbeiten (\S~22), Akkordarbeit (\S~23), Arbeiten unter
Tage (\S~24).
\end{loesung}

\subsection*{Lösung Aufgabe~6}
\begin{loesung}
\textbf{(a)} \emph{Während der Probezeit}: beide Seiten jederzeit, ohne Frist und ohne Grund
(\S~22 Abs.~1). \emph{Nach der Probezeit}: fristlos nur aus wichtigem Grund (außerordentlich,
innerhalb von zwei Wochen nach Bekanntwerden auszusprechen); ordentlich mit vier Wochen
Frist nur durch die/den Auszubildende:n (Aufgabe der Ausbildung / anderer Beruf) --- der
Betrieb hat dieses Recht nicht (\S~22 Abs.~2, 4).

\textbf{(b)} \emph{Einfaches Zeugnis}: Art, Dauer und Ziel der Ausbildung sowie erworbene
Fertigkeiten, Kenntnisse und Fähigkeiten. \emph{Qualifiziertes Zeugnis}: auf Verlangen
zusätzlich Angaben über \emph{Verhalten und Leistung}. Die elektronische Form ist
ausgeschlossen.

\textbf{(c)} Ein Betriebsrat wird ab in der Regel \emph{fünf} ständigen wahlberechtigten
Arbeitnehmern (davon drei wählbar) gewählt (\S~1 BetrVG); die Amtszeit beträgt
\emph{vier Jahre} (\S~21).

\textbf{(d)} Beispiel: Die JAV beantragt beim Betriebsrat Maßnahmen zugunsten der jungen
Beschäftigten (z.\,B.\ zu Berufsbildung und Übernahme) und wacht über die Einhaltung der für
Jugendliche und Auszubildende geltenden Vorschriften.
\end{loesung}

\end{document}
